\section{Code as Modality: por qué BLEU no basta}

El código es a la vez una token sequence y algo con semántica ejecutable. Dos cadenas muy distintas pueden implementar la misma función, y dos cadenas parecidas pueden diferir por un solo símbolo que provoca un error grave. Por eso, la evaluación de la code generation debe ejecutar la specification siempre que sea posible, en lugar de solo comparar el reference overlap.

\subsection{¿Es el código solo otra modality?}

Esta sección vuelve a las salidas que pueden verificarse automáticamente. El código puede seguir la token generation recipe, pero el parser, el runtime y los tests hacen que «correcto» sea más operable, y también más estricto, que la similitud en lenguaje natural. El texto y las imágenes vistos antes dependían sobre todo de evaluación humana o perceptual; con el código, por primera vez, la distribución de candidatos puede entrar en un ciclo cerrado de ejecución automática, clasificación de fallas y búsqueda.

\lecturefigure{slide-29-code-modality.jpg}{El código puede generarse con un modelo de lenguaje, pero su corrección está restringida por la semántica de ejecución.}{Video oficial 32:20.}

\begin{knowledgebox}{Lectura de la figura: la misma recipe, con un verifier adicional}
La tokenización del código y el autoregressive objective pueden reutilizarse directamente de GPT, pero el código cuenta con parser, type checker, tests y runtime. Estos proporcionan una retroalimentación automática más fuerte que la del lenguaje natural. Al mismo tiempo, los programas plantean problemas de security, resource, dependency y specification ambiguity; pasar un número limitado de tests es solo una de las evidencias.
\end{knowledgebox}

\teachervoice{El ponente pregunta «isn't code just another modality?» y de inmediato pasa a HumanEval. Nota de clase: la interfaz del modelo puede unificarse, pero la evaluation no admite atajos; la ventaja del código es que tiene un executable verifier, y la responsabilidad también es mayor.}

\subsection{HumanEval: problema, firma, docstring y tests}

\term{HumanEval} es un conjunto de Python function problems escritos a mano. Cada problema proporciona la function signature y el docstring, y el modelo genera el body; unas \term{unit tests} (pruebas unitarias) ocultas revisan varias entradas. \term{functional correctness} (corrección funcional) significa que el programa produce el comportamiento correcto en los specification cases, no que su texto se parezca al reference code.

\lecturefigure{slide-30-humaneval-dataset.jpg}{HumanEval evalúa el código generado con una function specification y hidden tests.}{Video oficial 33:33; Chen et al. 2021.}

\begin{knowledgebox}{Lectura de la figura: qué contratos incluye un problema}
La signature define la interfaz, el docstring da la specification en lenguaje natural, los examples pueden sugerir los casos límite y los tests definen la corrección observable. Si los tests no cubren entradas anómalas, la complejidad o los efectos secundarios, un programa erróneo todavía puede aprobarlos. La benchmark quality depende de la specification clarity y de la test coverage.
\end{knowledgebox}

\lecturefigure{slide-31-humaneval-example.jpg}{El HumanEval example muestra la relación entre el prompt y la implementation candidata.}{Video oficial 35:15.}

\begin{knowledgebox}{Lectura de la figura: primero razonar a mano, luego ejecutar los tests}
Lea el nombre de la función, los parámetros y el docstring, y enumere los normal, boundary, empty, duplicate e invalid cases; después revise el código candidato. El modelo puede escribir una implementación sintácticamente correcta y aparentemente razonable que, sin embargo, falla por off-by-one, mutation o límites de tipo. La ejecución es una verificación necesaria, pero un secure sandbox y un timeout también son indispensables.
\end{knowledgebox}

\begin{lstlisting}[caption={Forma simplificada de un functional check al estilo HumanEval}]
def candidate(values):
    # generated implementation
    ...

assert candidate([1, 2, 3]) == expected_normal
assert candidate([]) == expected_empty
assert candidate(boundary_case) == expected_boundary
\end{lstlisting}

\subsection{Pass@k: al menos una muestra correcta en la distribución}

Si se muestrean $k$ candidatos independientes del modelo, \term{pass@k} indica la probabilidad de que al menos uno pase los tests. Si para cada problema se muestrean $n$ programas, de los cuales $c$ son correctos, el artículo de Codex usa el estimador
\[
\widehat{\operatorname{pass@}k}
=1-\frac{\binom{n-c}{k}}{\binom{n}{k}},
\]
y cuando $n-c<k$ el resultado es 1. El numerador es el número de combinaciones de $k$ muestras elegidas entre las incorrectas, y el denominador es el número de todos los subconjuntos de tamaño $k$; su cociente es la probabilidad de que «todas las elegidas sean incorrectas».

\lecturefigure{slide-32-pass-at-k.jpg}{Pass@k evalúa la probabilidad de encontrar al menos un programa funcionalmente correcto en varios muestreos.}{Video oficial 36:00; Chen et al. 2021.}

\begin{knowledgebox}{Lectura de la figura: pass@1 y pass@100 responden a preguntas de producto distintas}
Pass@1 se aproxima a la experiencia de un usuario que solo ve una completion; un large-k se acerca más a un sistema que puede generar muchos candidatos y elegir con tests/reranker. Un pass@100 alto con un pass@1 bajo indica que los programas correctos existen en la distribución, pero que el decoding por defecto rara vez los alcanza. El valor de despliegue depende de si realmente existe un verifier, del costo de muestreo y de la latencia de selección.
\end{knowledgebox}

\begin{warningbox}{Pass@k no es la proporción de «el modelo sabe escribir software»}
Los problemas de HumanEval son cortos, con specification clara y tests automatizados. El software real incluye dependencias entre archivos, requisitos ambiguos, security, performance, maintenance y review. El pass@k de un benchmark es evidencia de una capacidad acotada, no de autonomous software-engineering reliability.
\end{warningbox}

\subsection{Codex training snapshot}

La slide de la clase registra que los datos de Codex de entonces eran 159 GB de code recopilados de 54 millones de repositories, y que se hizo fine-tuning a partir de GPT-3 models de distintos tamaños; el tokenizer añadió soporte para spaces. Esta es una descripción histórica de la época del artículo y no representa los datos de productos posteriores ni la política actual.

\lecturefigure{slide-33-codex-training-details.jpg}{Codex parte de GPT-3, aplica fine-tuning con un code corpus a gran escala y ajusta el tokenizer.}{Video oficial 36:59; Chen et al. 2021.}

\begin{knowledgebox}{Lectura de la figura: además del volumen de datos, hay una representation choice}
El repository count no se traduce directamente en examples independientes y de alta calidad; se requieren deduplicación, revisión de licencias, filtrado de seguridad y una train-test contamination audit. Los spaces tienen semántica en la indentación de Python, y el cambio en el tokenizer puede reducir la longitud de las secuencias y mejorar el aprendizaje de la estructura. El fine-tuning reutiliza las capacidades de lenguaje natural y también conecta el prompt en lenguaje natural con la code distribution.
\end{knowledgebox}

\subsection{Resumen de la sección}

La code generation hereda la language-model recipe, pero gracias a la executable semantics obtiene un verifier más fuerte. HumanEval y pass@k elevan el criterio de «parece código» a «al menos un candidato pasa los tests», aunque siguen cubriendo solo una pequeña parte del software real.
